\documentclass[11pt,a4paper]{article}
\usepackage[T1]{fontenc}
\usepackage[utf8]{inputenc}
\usepackage[slovene]{babel}
\usepackage{amsmath,amssymb}
\usepackage{graphicx}
\usepackage{booktabs}
\usepackage{xcolor}
\usepackage{listings}
\usepackage[margin=2.7cm]{geometry}
\usepackage[hidelinks]{hyperref}
\usepackage{placeins}

\renewcommand{\topfraction}{0.9}
\renewcommand{\bottomfraction}{0.8}
\renewcommand{\textfraction}{0.07}
\renewcommand{\floatpagefraction}{0.75}
 
\newcommand{\slikaali}[2]{\IfFileExists{#1}{\includegraphics[width=#2\linewidth]{#1}}%
  {\dopolni{sliko \texttt{\detokenize{#1}} ustvari \texttt{vizualizacije.py}}}}
\lstset{language=Python,basicstyle=\ttfamily\small,frame=single,breaklines=true,
        showstringspaces=false}
 
\title{Lastne vrednosti simetrične matrike\\ s QR iteracijo z enojnim premikom}
\author{Urban Vesel}
\date{10. avgust 2026}
 
\begin{document}
\maketitle
 
\section{Opis problema}
 
Za simetrično matriko $A \in \mathbb{R}^{n\times n}$ iščemo vse njene lastne
vrednosti in po želji pripadajoče lastne vektorje. Po spektralnem izreku ima
vsaka taka matrika realne lastne vrednosti in ortonormirano bazo lastnih
vektorjev, torej razcep $A = V\Lambda V^{\mathsf T}$ z ortogonalno $V$ in
diagonalno $\Lambda$.
 
Lastnih vrednosti za $n > 4$ ne moremo poiskati z zaključeno formulo, zato jih
računamo iterativno. Uporabimo QR iteracijo s premikom: matriko postopoma
prevajamo v diagonalno obliko z zaporedjem ortogonalnih podobnostnih
transformacij, ki spekter ohranjajo, obdiagonalne elemente pa potiskajo
proti nič.
 
\section{Opis rešitve}
 
\subsection{Osnovna ideja}
 
Postopek ima dva koraka. Najprej matriko s Householderjevimi zrcaljenji
prevedemo na simetrično \emph{tridiagonalno} obliko $T$. To je končen
postopek, ki matriko poenostavi, spektra pa ne spremeni. Nato tridiagonalno
matriko iterativno diagonaliziramo: v vsakem koraku izberemo premik $\mu$,
izračunamo QR razcep premaknjene matrike in faktorja zmnožimo v obratnem
vrstnem redu,
\begin{equation}
  T - \mu I = QR, \qquad T \leftarrow RQ + \mu I .
  \label{eq:korak}
\end{equation}
Ker je $RQ + \mu I = Q^{\mathsf T} T Q$, je vsak korak ortogonalna podobnostna
transformacija in lastne vrednosti se ohranijo natančno; premik vpliva samo na to,
kako hitro se posamezen obdiagonalni element približa ničli. Ko zadnji
obdiagonalni element postane zanemarljiv, se zadnja lastna vrednost odcepi
(deflacija). Postopek nadaljujemo na manjši matriki.
 
Tridiagonalna oblika se pri koraku \eqref{eq:korak} ohrani. Matrika $R$ ima nad
diagonalo le dva neničelna pasova, $Q$ je zgornja Hessenbergova, zato je tak
tudi produkt $RQ$; ker je hkrati simetričen, je nujno spet tridiagonalen. Zaradi
tega lahko ta cel korak izvedemo neposredno na diagonali in obdiagonali, z
zaporedjem Givensovih rotacij in brez hranjenja polne matrike.
 
\subsection{Zakaj je enojni premik dobra izbira}
 
Naloga predpisuje premik $\mu = t_{nn}$, torej zadnji diagonalni element
aktivnega bloka. Iz $Q^{\mathsf T}(T-\mu I) = R$ za
zadnjo vrstico dobimo
\begin{equation}
  q_n = r_{nn}\,(T-\mu I)^{-1} e_n ,
\end{equation}
kar pomeni, da je zadnji stolpec matrike $Q$ natanko rezultat enega koraka
inverzne iteracije s premikom $\mu$. Ker je $\mu = e_n^{\mathsf T} T e_n$
Rayleighov kvocient trenutnega približka, je celoten postopek Rayleighova
kvocientna iteracija, zapisana v bazi, ki jo sproti vrtimo. Pri simetričnih
matrikah je Rayleighov kvocient od lastne vrednosti oddaljen kvadratično glede
na kotno napako približka, zato je konvergenca \emph{kubična}: če je napaka v
enem koraku $\varepsilon$, je v naslednjem reda $\varepsilon^3$. To se lepo vidi
v zaporedju obdiagonalnega elementa med iteracijo,
\[
  6{,}5\cdot 10^{-1} \;\to\;
  8{,}6\cdot 10^{-2} \;\to\;
  3{,}8\cdot 10^{-4} \;\to\;
  3{,}8\cdot 10^{-10} \;\to\;
  7{,}7\cdot 10^{-28},
\]
kjer je vsak člen približno tretja potenca prejšnjega (slika~\ref{fig:konv}).

Iz zadnjih treh členov zaporedja dobimo oceno reda konvergence $p = \frac{\log(e_{k+1}/e_k)}{\log(e_k/e_{k-1})} = 2{,}95$, 
kar se dobro ujema z teoretičnim pričakovanjem $p = 3$.
 
\subsection{Deflacija in zanesljivost}
 
Obdiagonalni element zanemarimo, ko velja
$|e_i| \le \varepsilon_{\mathrm{stroj}}\,(|d_i| + |d_{i+1}|)$. Postavitev takega
elementa na nič je motnja matrike z normo $|e_i|$, po Weylovem izreku se
lastne vrednosti pri motnji premaknejo kvečjemu za njeno normo. Kriterij je
torej povratno stabilen in vezan na velikost matrike, ne na absolutno vrednost.
 
Enojni premik pa ima slabost, da lahko obtiči. Za matriko
$\bigl[\begin{smallmatrix}0&1\\1&0\end{smallmatrix}\bigr]$ je $\mu = 0$ in korak
matrike sploh ne spremeni, ker je Rayleighov kvocient točno na sredini med
lastnima vrednostma. Zato iteracija spremlja padanje zadnjega obdiagonalnega
elementa in ob zastoju preklopi na Wilkinsonov premik, ki konvergira vedno. 
Pri naših poskusih se je varovalo sprožilo v 3 od 1284 korakov (0,2 \%), 
na naključnih matrikah torej skoraj nikoli. Izjema je antidiagonalna $2\times 2$ matrika, 
pri kateri enojni premik dokazano obtiči: tam varovalo prevzame in iteracijo odblokira.
 
\subsection{Zahtevnost}
 
Redukcija na tridiagonalno obliko je najdražji del. Stane $\mathcal{O}(n^3)$
operacij. Posamezen korak iteracije pa se zaradi tridiagonalne strukture izvede
v $\mathcal{O}(n)$ (oziroma $\mathcal{O}(n^2)$, če hkrati posodabljamo lastne
vektorje), prostor pa je zgolj $\mathcal{O}(n)$, saj polne matrike nikoli ne
hranimo. Ker zaradi kubične konvergence na lastno vrednost zadostuje le nekaj
korakov (izmerili smo približno $2{,}6$ do $3$ korake), je iteracijski del skupaj
$\mathcal{O}(n^2)$.
 
\section{Rezultati}
 
\begin{figure}[htbp]\centering
  \slikaali{slike/konvergenca.png}{0.6}
  \caption{Padanje zadnjega obdiagonalnega elementa po korakih iteracije za
  naključno simetrično matriko velikosti $12\times 12$. Navpična os je
  logaritemska; padec z $10^{-4}$ na $10^{-10}$ in nato na $10^{-28}$ v dveh
  korakih je značilen za kubično konvergenco.}
  \label{fig:konv}
\end{figure}
 
Natančnost smo preverili na naključnih simetričnih matricah in jo primerjali z
vgrajeno funkcijo \texttt{numpy.linalg.eigvalsh}:
 
\begin{table}[htbp]\centering
\begin{tabular}{rccc}
\toprule
$n$ & $\max|\Delta\lambda|$ & $\|V^{\mathsf T}V - I\|_{\max}$
    & $\|AV - V\Lambda\|_{\max}$ \\ \midrule
 10 & $2{,}2\cdot10^{-15}$ & $1{,}2\cdot10^{-15}$ & $2{,}0\cdot10^{-15}$ \\
 40 & $2{,}0\cdot10^{-14}$ & $6{,}0\cdot10^{-15}$ & $9{,}8\cdot10^{-15}$ \\
100 & $7{,}8\cdot10^{-14}$ & $9{,}1\cdot10^{-15}$ & $2{,}1\cdot10^{-14}$ \\
200 & $2{,}3\cdot10^{-13}$ & $2{,}2\cdot10^{-14}$ & $5{,}8\cdot10^{-14}$ \\
\bottomrule
\end{tabular}
\caption{Odstopanje lastnih vrednosti od oraklja, ortogonalnost izračunane
baze in ostanek lastnega problema. Napaka narašča počasi z velikostjo matrike
in ostaja blizu strojne natančnosti.}
\end{table}
 
\begin{figure}[htbp]\centering
  \slikaali{slike/casi.png}{0.6}
  \caption{Čas izračuna v odvisnosti od velikosti matrike. Izmerjeni naklon je
  približno $1{,}84$, čeprav je metoda kot celota $\mathcal{O}(n^3)$: pri teh velikostih redukcija, 
  izvedena z matričnimi operacijami knjižnice numpy, 
  še ne prevlada, prevladuje pa iteracijski del, ki je v Pythonu izveden kot navadna zanka in je
  $\mathcal{O}(n^2)$. Vgrajena funkcija je zato približno stokrat hitrejša ob enaki asimptotiki.}
  \label{fig:casi}
\end{figure}


Hitrost smo primerjali z isto vgrajeno funkcijo (slika~\ref{fig:casi}).
Pri $n = 200$ je naša implementacija približno stokrat počasnejša, izmerjeni naklon pa je $1{,}84$.


Meritve so bile opravljene z numpy 2.5.1; 
ker so razlike v zadnjih bitih odvisne od izvedbe BLAS, 
se lahko števila korakov in napake na drugi napravi razlikujejo 
za nekaj odstotkov.
 
\section{Primer uporabe: spektralno gručenje}
 
Množico točk lahko razdelimo v gruče s pomočjo lastnih vektorjev. Točke
povežemo v utežen graf, kjer je utež med dvema točkama
$w_{ij} = \exp(-\|x_i - x_j\|^2 / 2\sigma^2)$, torej blizu ena za bližnji
točki in skoraj nič za oddaljeni. Laplaceova matrika grafa je $L = D - W$, kjer
je $D$ diagonalna matrika vsot vrstic. Najmanjša lastna vrednost take matrike je
vedno nič, pripadajoči lastni vektor pa je konstanten. Predznak komponent
drugega lastnega vektorja (Fiedlerjevega) pa točke razdeli na dva dela tako, da
je vsota uteži prerezanih povezav majhna.


Na dveh Gaussovih oblakih razrez pravilno razvrsti vse točke, 
med drugo in tretjo lastno vrednostjo pa je velika vrzel (vidno na sliki~\ref{fig:gruce}).
 
\vspace{4em}

\begin{lstlisting}
import numpy as np
from dn1_qr_premik import (EnojniPremik, eigen,
                           laplaceova_matrika, podobnostna_matrika)
 
L = laplaceova_matrika(podobnostna_matrika(X, sigma=0.8))
lastne, V = eigen(L, EnojniPremik(), vektorji=True)
gruca = V[:, 1] > 0        # razrez po Fiedlerjevem vektorju
\end{lstlisting}
 
\begin{figure}[htbp]\centering
  \slikaali{slike/gruce.png}{0.46}\hfill
  \slikaali{slike/spekter.png}{0.46}
  \caption{Levo: 130 točk iz dveh Gaussovih oblakov, pobarvanih po predznaku
  Fiedlerjevega vektorja; razrez pravilno razvrsti vseh 130 točk. Desno:
  spodnji del spektra Laplaceove matrike. Prva lastna vrednost je
  $8{,}4\cdot 10^{-15}$, torej numerično nič, med drugo in tretjo pa je velika
  vrzel ($\lambda_3/\lambda_2 \approx 45$), kar je znak, da sta gruči res dve.}
  \label{fig:gruce}
\end{figure}

\FloatBarrier
 
\section{Zaključek}
 
Implementirali smo iskanje lastnih vrednosti in lastnih vektorjev simetrične
matrike s QR iteracijo z enojnim premikom, pri čemer je posamezen korak izveden
z Givensovimi rotacijami neposredno na diagonali in obdiagonali. Rezultati se z
vgrajeno funkcijo ujemajo do reda $10^{-13}$, izračunana baza pa ostane
ortogonalna do strojne natančnosti. Kubična konvergenca se v poskusih jasno
pokaže: na lastno vrednost v povprečju zadostujeta dva do trije koraki.
 
Glavna omejitev je hitrost. Ker je iteracijska zanka napisana v Pythonu, je
implementacija pri $n = 200$ približno stokrat počasnejša od vgrajene funkcije,
čeprav je asimptotika enaka. Naravna izboljšava bi bila prevod te zanke v prevedeno
kodo ali uporaba implicitnega koraka, ki premaknjene matrike sploh ne
izračuna eksplicitno.

\FloatBarrier
 
\begin{thebibliography}{9}
\bibitem{gvl} G.~H. Golub, C.~F. Van Loan, \emph{Matrix Computations},
  4.~izdaja, Johns Hopkins University Press, 2013, razdelek~8.3.
\bibitem{parlett} B.~N. Parlett, \emph{The Symmetric Eigenvalue Problem},
  SIAM, 1998.
\bibitem{skripta} M.~Vuk, \emph{Numerična matematika v programskem jeziku
  Julia}, 2025, poglavje~8 in naloga~18.1.9.
\end{thebibliography}
 
\end{document}
